\section{Constant-predictor baseline}
\label{app:constant}

The chance level of $\pi/2 + \nicefrac{2}{\pi} \approx 2.207$\,rad assumes orientations drawn
uniformly from $SO(3)$, which SPEED+ poses are not. This appendix gives the empirical
counterpart quoted in Section~\ref{sec:setup}: the score of a single fixed pose. It anchors
``near chance'' to a predictor one could actually build, rather than to an assumption about
the distribution.

The constant is fitted on the 28{,}000 training labels only. Fitting it on a test domain
would make it an oracle rather than a baseline, so neither real domain enters the estimate.
For rotation we take the principal eigenvector of $M = \sum_i \mathbf{q}_i
\mathbf{q}_i^{\top}$, which minimises chordal distance and is unaffected by the quaternion
double cover. We then refine it with five Karcher iterations on the sphere: sign-align each
$\mathbf{q}_i$ to the current mean, log-map into the tangent space, average, and exp-map back.
The $SO(3)$ geodesic distance is twice the $S^3$ one, so iterating on the sphere targets the
same minimiser as our reported metric. For translation we take the arithmetic mean of the
training translations. This gives $\mathbf{q} = (-0.198, 0.034, -0.019, -0.979)$ up to sign
and $\mathbf{t} = (-0.0001, 0.0026, 6.120)$\,m.

Two caveats belong with the numbers. First, the Karcher steps do not contract sharply here,
because the training orientations cover most of the sphere, which leaves the Fr\'echet
functional close to flat and the mean weakly identified. This does not affect the result: the
chordal and
Karcher estimates lie $0.022$\,rad apart and agree to $0.001$ on every score in
Table~\ref{tab:constant}. Second, both estimators minimise \emph{squared} error, whereas we
report mean unsquared error, whose minimisers are the corresponding geometric medians. The
scores below are therefore achievable rather than optimal, and they are an upper bound on the
error of the best constant predictor. We do not claim optimality, and no conclusion in the
paper needs it.

\begin{table}[ht]
  \caption{A constant pose fitted on the training split, evaluated on each test
  domain. $\Etrans$ range-normalised, $\Erot$ in radians. Skill is
  $2.207 - \Erot$. Reported for the Karcher-refined estimate.}
  \label{tab:constant}
  \centering
  \small
  \begin{tabular}{l c cc c}
    \toprule
    Domain & $n$ & $\Etrans$ & $\Erot$ & Skill \\
    \midrule
    synthetic (val) & 2{,}000 & 0.415 & 2.176 & 0.031 \\
    lightbox & 6{,}740 & 0.367 & 2.203 & 0.005 \\
    sunlamp & 2{,}791 & 0.372 & 2.199 & 0.009 \\
    \bottomrule
  \end{tabular}
\end{table}

The empirical no-skill rotation level therefore sits at most $0.031$\,rad from the uniform
one on all three domains, and we use $2.207$ in the body. Every trained run behind
Table~\ref{tab:results} keeps some skill over it, even the weakest, seed 42 of the frozen
un-augmented arm on sunlamp, which scores $2.048$\,rad and keeps $0.15$\,rad.

Translation is less comfortable, and it is where the second caveat matters most. On sunlamp
the constant scores $0.372$, which two of the three augmented layer-4 seeds ($0.405$, $0.399$)
do not beat, while a constant fitted directly to the unsquared metric does better on
translation, scoring
$0.330$ on sunlamp, and matches this one on rotation to within $0.001$. We quote the
Karcher-refined constant throughout, as the workshop version did, because no conclusion
changes with the tighter one: the same two seeds fail to beat it, and every other comparison
the paper draws with the constant still holds.
